% !PRINTED 292
\weberpage{292}

According to (3) we can derive from this the composition series of the two groups $P/Q$
and $P/R$, which we shall indicate thus:
\weq{5}{P/Q,\ Q'/Q,\ Q''/Q,\ \ldots,\ 1}
\weq{6}{P/R,\ Q'/R,\ Q''/R,\ \ldots,\ Q/R \ldots}
Now the index of the divisor $Q'/Q$ of $P/Q$ is equal to the index of the divisor $Q'$
of $P$, and hence also equal to the index of the divisor $Q'/R$ of $P/R$ (\S.\,6, 1.), and
the same holds of the following terms. We therefore state the theorem:

\noindent 3. \gesperrt{If $Q$ and $R$ are normal divisors of $P$, and $R$ is a
divisor of $Q$, then the index series of $P/Q$ is a part of the index series of
$P/R$..} % [sic] double period as printed

If the indices $j_1, j_2, \ldots, j_\nu$ are all primes, then the solution
of the resolvent $\chi(y) = 0$ is reducible to the solution of a chain of cyclic
equations of degrees $j_1, j_2, \ldots, j_\nu$; this resolvent is
therefore metacyclic in the sense in which we settled this concept in the seventeenth chapter
of the first volume, according to which the metacyclic equations are
identical with those otherwise called algebraically solvable.

An irreducible equation $f(x) = 0$ is metacyclic if the index series of its
group consists entirely of primes. At the place mentioned we investigated the
conditions for metacyclic equations of prime degree, and
must now carry these considerations through for the general case, in which
the degree of the equation is arbitrarily composite.

\begin{center}
\S.\,74.

\medskip

\gesperrt{Metacyclic Equations.}
\end{center}

\bigskip

If $f(x) = 0$ is an irreducible equation of degree $m$ and $P$ is its
\gesperrt{Galois} group, then $P$ is transitive. Let a composition and index series
for $P$ be
\weq{1}{
\begin{gathered}
P,\ P_1,\ P_2,\ \ldots,\ P_{\mu-1},\ 1 \\
j_1,\ j_2,\ \ldots,\ j_{\mu-1},\ j_\mu.
\end{gathered}
}
Since the last group $1$ is intransitive, an intransitive group must at some point
occur in the composition series; let $P_\lambda$ therefore be the first
intransitive group of the sequence (1). Then
%CONT>>
